% 00-map.tex — bản đồ Phần E: có gì, mức đầu tư, thứ tự đọc.
% Ngày tạo: 2026-09-03 | Sửa gần nhất: 2026-09-03 | Ôn gần nhất: chưa
% Dependency: ../00-map.tex | Mức độ: điều hướng
\documentclass[../deep_learning.tex]{subfiles}
\begin{document}
\begin{table}[H]\centering\small
\caption{Bản đồ Phần E: các chương, nội dung, và mức đầu tư đề xuất.}
\begin{tabularx}{\linewidth}{@{}L{0.8cm}L{4.0cm}YL{2.6cm}@{}}
\toprule
\textbf{Ch.} & \textbf{Thư mục} & \textbf{Nội dung} & \textbf{Mức}\\
\midrule
23 & \code{23-doc-paper} & Đọc ba lượt; đọc giới hạn trước kết quả; ghi chú phải nêu chỗ mình không tin & \textbf{cốt lõi}\\
24 & \code{24-thi-nghiem-va-ablation} & Thiết kế so sánh, ablation, thống kê tối thiểu, phân tích lỗi, báo cáo & \textbf{cốt lõi}\\
25 & \code{25-do-tin-cay-va-van-hanh} & Trực quan hoá và giới hạn của nó, hiệu chỉnh, bất định, vận hành & cốt lõi\\
26 & \code{26-lo-trinh-hoc} & Ba vòng dự án, gợi ý đề tài, và thứ tự học đề xuất cho cả cây & \textbf{cốt lõi}\\
\bottomrule
\end{tabularx}
\end{table}

\noindent Chương 23 và 24 nên đọc song song với phần C chứ không đợi tới cuối, vì chúng thay đổi cách bạn đọc mọi paper trong phần C. Chương 26 chứa thứ tự học đề xuất cho toàn cây.
\end{document}
